\section{Related Work}

\textbf{Foundations of Subject-Driven Generation.} 
The rapid evolution of text-to-image diffusion models~\cite{rombach2022high} has enabled profound new capabilities in personalized image generation. Foundational techniques like DreamBooth~\cite{ruiz2023dreambooth} and Textual Inversion~\cite{gal2022image} first demonstrated that a pre-trained model could learn novel subjects using only a few reference images. To improve efficiency and spatial control, subsequent works introduced parameter-efficient fine-tuning like LoRA~\cite{hu2021lora}, structural conditioning such as ControlNet~\cite{zhang2023adding}, and direct image-prompt adapters like IP-Adapter~\cite{ye2023ip}. While these methods achieve remarkable single-subject fidelity, they struggle with severe attribute binding and identity bleeding when tasked with generating multiple distinct subjects simultaneously.

\textbf{Multi-Subject Personalization: Solving "Who is Who".} 
Under DiT \cite{peebles2023scalable} and FLUX \cite{flux2024} backbones, early multi-subject research primarily addressed a subject-centric question: how can multiple referenced identities remain distinct under shared generation dynamics? XVerse \cite{xverse2024} is a representative method in this direction. It converts reference images into token-specific text-stream modulation offsets, injecting subject-specific control along the text-conditioning pathway rather than directly perturbing image latents. This non-invasive design improves identity controllability while avoiding direct corruption of the shared image representation. Similarly, MOSAIC \cite{mosaic2024} tackles this problem from a representation-centric perspective. By enforcing semantic correspondence and feature disentanglement across references, it explicitly separates subjects in the feature space, reducing identity blending. Taken together, these methods move personalization beyond simple fidelity preservation toward explicit subject separation, establishing the first layer of controllable multi-subject generation.

\textbf{Layout-to-Region Binding: Solving "Where Each Subject Goes".} 
Once subject identity is preserved, the next challenge is spatial grounding. This motivates a second line of work on layout-to-region binding, aiming to align references, prompts, and spatial regions. ContextGen \cite{contextgen2024} introduces Contextual Layout Anchoring (CLA) and Identity Consistency Attention (ICA) to stabilize the correspondence between instances and target positions. AnyMS \cite{anyms2024} addresses the same challenge in a training-free manner by decoupling attention at global and local levels to jointly maintain prompt alignment and layout fidelity. 3DIS-FLUX \cite{3disflux2024} approaches the problem through depth-guided scene construction followed by FLUX-based rendering, utilizing depth as a compact structural prior. However, when the depth map underspecifies hidden geometry or front-back ordering, occlusion relations may still collapse. In essence, layout binding methods solve the "where" problem but only partially address the physically harder question of what should remain visible under overlap.

\textbf{Attention Leakage and Entanglement in DiT/FLUX.} 
To understand multi-subject failures, it is crucial to examine the mechanics of attention leakage. In DiT-style generators, global attention is both a strength and a liability. While it enables rich long-range interaction, it also allows identity cues, attributes, colors, and poses from one subject to leak into another. This issue is especially pronounced in FLUX \cite{flux2024}, where the later single-stream blocks flatten text and image modalities into a tightly fused sequence, potentially losing the structural inductive bias for distinct instances. From this perspective, semantic entanglement is not merely an identity-preservation failure, but an \textit{attention-routing failure}. Existing methods attempt to mitigate this through various strategies—XVerse via text-stream modulation, MOSAIC via semantic disentanglement, and AnyMS via decoupled attention flows. This progression suggests that multi-subject generation is fundamentally constrained by how well the model can route attention without cross-instance contamination.

\textbf{Interaction and Occlusion Reasoning.} 
The final and most challenging layer is physical interaction and occlusion reasoning. At this stage, the problem shifts to who occludes whom, whether partially hidden instances remain geometrically plausible, and whether the model respects front-back ordering under overlap. SeeThrough3D \cite{seethrough3d2024} explicitly models which object should appear in front and how partial visibility should be rendered consistently, connecting naturally to the notion of amodal completion. LayerBind \cite{layerbind2024} leverages the coarse-to-fine denoising dynamics of DiTs, establishing layout and visible order at early steps through layered initialization. Nevertheless, deeply entangled interactions—such as hugging, grappling, or interleaved articulated bodies—remain exceptionally difficult. The challenge lies not only in visibility ordering but also in fine-grained contact geometry, mutual deformation, and physically consistent completion under occlusion. Our work explicitly targets this gap, providing a rigorous benchmark and novel evaluation metrics to measure where current models fail in physically grounded interaction reasoning, perfectly aligning with the P13N workshop's focus on personalization quality and multi-subject composition.
